\documentclass[11pt,letterpaper]{article}
\usepackage{jahnel}
\runningtitle{Occupied passage through oscillating corridors}
\hypersetup{pdftitle={Occupied passage through oscillating corridors},pdfauthor={}}
\title{Occupied passage through oscillating corridors}
\author{}
\date{November 3, 2026}

\begin{document}
\maketitle
\begin{abstract}
We estimate occupied passage through the oscillating configuration used in the preceding environmental lower bound. Conditional on the selected lines, we restore the remaining Poisson line process and place independent Poisson points on every line. Passage through a tube of fixed width following the selected graph has probability between $\exp\{-Ck(\log k)^2\}$ and $\exp\{-ck(\log k)^2\}$, where the total horizontal length is proportional to $k^2$. The upper bound allows every occupied route inside the tube. Its proof bounds the total contribution of one line to a sequence of separating rectangles and then applies the Poisson exponential formula. The lower bound inserts a bundle of nearly parallel lines along each supporting segment and simultaneously preserves the original large count of full connected runs. Thus, sampling the Cox points adds a substantial connection cost, but the complete joint cost remains subexponential in the total length. The estimate concerns a prescribed tube and leaves the isotropic subcritical problem open.
\end{abstract}

\section{The occupied-passage estimate}

The preceding note constructs oscillating corridors with many full connected runs at probability cost $\exp\{-O(k\log k)\}$ over a total of $k^2$ corridor widths~\cite{OSC26}. That calculation concerns the line environment. We now test whether occupied passage through the same geometry has a substantially larger cost after the Cox points are sampled. The estimate below retains every remaining line and permits all routes inside a tube following the selected graph.

Let $\mu,\lambda,K>0$. The Poisson line process $\cL$ has intensity $\nu_\mu(\dd\theta\,\dd s)=(\mu/\pi)\dd\theta\dd s$ on lines $\ell(\theta,s)=\{x\in\R^2:x\cdot(\cos\theta,\sin\theta)=s\}$, where $0\le\theta<\pi$ and $s\in\R$. Conditional on $\cL$, place an independent Poisson process of intensity $\lambda$ per unit arclength on every line. Let $X$ be their union and $\Xi=\bigcup_{z\in X}B(z,1)$, where $B(z,a)$ is the open Euclidean disk of radius $a$ centred at $z$. Write $\P$ and $\E$ for probability and expectation, and $\mathcal H^1$ for arclength.

We retain the parameters and selection sets of~\cite{OSC26}. Fix an integer $t\ge\max(8,4\mu)$, put $w=(t/\mu)^{3/2}$, and assume $w/t\ge4096\e^{4\mu}(K+4)$. For positive integers $k$ divisible by $8t$, define
\begin{equation}\label{eq:parameters}
 M=k^2,\qquad L=kw,\qquad R=kL,\qquad b=\frac1{4t},\qquad
 H=\frac{bL}{2},\qquad x_i=iL\quad(0\le i\le k).
\end{equation}
Let $\zeta_0=0$ and $\zeta_i=(-1)^{i+1}H$ for $i\ge1$. The continuous piecewise affine function $F_0$ joins the points $(x_i,\zeta_i)$. Set $c_i=(i-1/2)L$, $d_i=(\zeta_{i-1}+\zeta_i)/2$, and $\beta_i=(\zeta_i-\zeta_{i-1})/L$. Thus, the first slope is $b/2$ and the subsequent slopes alternate between $-b$ and $b$.

Put $\varepsilon=1/100$ and $N=\lfloor H/64\rfloor$. For a graph line $h(x)=a+vx$, define
\begin{equation}\label{eq:selection}
\begin{aligned}
 \mathscr S_i&=\{h:|h(c_i)-d_i|\le\varepsilon,
                           \ |v-\beta_i|\le\varepsilon/L\},\quad 1\le i\le k,\\
 \mathscr B_i&=\{h:|h(c_i)-d_i|\le\varepsilon,
        \ |v-\beta_i+\operatorname{sign}(\beta_i)/L|\le\varepsilon/L\},
                                                        \quad 1\le i\le k,\\
 \mathscr H_j&=\{h:|h(0)-16j|\le\varepsilon,
                           \ |v|\le\varepsilon/(kL)\},\quad 1\le j\le N.
\end{aligned}
\end{equation}
Let $\mathscr U_k$ be their union, $E_k$ the event that each set contains exactly one line, and $\mathcal G_k$ the information supplied by $\cL$ inside $\mathscr U_k$. On $E_k$, denote the selected lines by $\sigma_i\in\mathscr S_i$, $\ell_i\in\mathscr B_i$, and $\gamma_j\in\mathscr H_j$, and their set by $\mathcal D_k$. The selection result in~\cite[Proposition~2.1]{OSC26} states that these sets are pairwise disjoint for sufficiently large $k$ and, with $\rho=w/(512t)$,
\begin{equation}\label{eq:selection-cost}
 \log\P(E_k)=-(2+2\rho)k\log k+O_{\mu,t}(k).
\end{equation}
This statement concerns the lines alone, so the additional point processes do not change it.

Put $a_0=12\varepsilon t$ and write $g_i$ for the affine function of $\sigma_i$. For $1\le i<k$, let $y_i$ be the horizontal coordinate of $\sigma_i\cap\sigma_{i+1}$, and set $y_0=0$, $y_k=R$. For sufficiently large $k$, $|y_i-x_i|\le12\varepsilon t$: the two heights at $x_i$ differ by at most $3\varepsilon$, and the slopes differ by at least $b$. The junctions are therefore ordered. Define the continuous graph $g$ by $g=g_i$ on $[y_{i-1},y_i]$.

The tube and its occupied-passage event are
\begin{equation}\label{eq:tube}
\begin{aligned}
 \mathcal T_k&=\{(x,y):1\le x\le R-1,\ |y-g(x)|\le2\},\\
 C_k&=\{\Xi\cap\mathcal T_k\text{ contains a continuous path from }x=1
                                      \text{ to }x=R-1\}.
\end{aligned}
\end{equation}
Centers outside the tube may contribute disks, and no remaining line is deleted. The curve itself must stay inside the tube. We also retain the environmental event that motivated the construction, now evaluated after every additional line has been restored.

For a base line $\ell$, split the other lines according to whether their normal angle is smaller or larger than that of $\ell$. Let $U_\ell^-$ and $U_\ell^+$ be the corresponding unions of intervals on $\ell$ at distance at most $2$ from the generating lines. A run is a full connected component of $U_\ell^-\cap U_\ell^+$. Let $I_i$ be the segment of $\ell_i$ with horizontal coordinates in $[y_{i-1},y_i]$, let $S_k$ count, with multiplicity over $i$, the full runs of arclength at least $t$ meeting $I_i$, and put $F_k=\{S_k\ge(K+4)k^2\}$. Components extending beyond a segment are included. Set the events involving the selected graph to be empty off $E_k$.

All probabilities conditional on $\mathcal G_k$ below are asserted almost surely on $E_k$. Constants $c,C>0$ and $k_0$ may depend on the fixed parameters $\mu,t,K,\lambda$.

The theorem gives matching orders for occupied passage, including passage that coexists with the original large full-run count. Its joint estimate also includes selection of the original configuration.

\begin{theorem}\label{thm:passage}
Constants $c,C,k_0$ can be chosen so that, for every $k\ge k_0$ divisible by $8t$,
\begin{equation}\label{eq:main}
\begin{aligned}
 \P(F_k\cap C_k\mid\mathcal G_k)
 &\ge\exp\{-Ck(\log k)^2\}\quad\text{on }E_k,\\
 \P(C_k\mid\mathcal G_k)
 &\le\exp\{-ck(\log k)^2\}\quad\text{on }E_k,\\
 \exp\{-Ck(\log k)^2\}
 &\le\P(E_k\cap F_k\cap C_k)\le\exp\{-ck(\log k)^2\}.
\end{aligned}
\end{equation}
\end{theorem}

In terms of $M$, these probabilities have logarithmic order $-\sqrt M(\log M)^2$. The estimate holds for every fixed $\lambda>0$. In particular, even at arbitrarily small point intensity, the joint probability cannot have an upper bound $A\exp\{-aM\}$ with finite $A$ and positive $a$. This conclusion concerns the stated passage event and does not establish percolation.

\section{Proof of the occupied-passage estimate}

The upper and lower bounds use different features of the oscillation. For the upper bound, a line can contribute only order $k$ arclength to a sequence of rectangles that every passage must meet. For the lower bound, order $\log k$ additional lines near each supporting segment suffice to carry an occupied path with conditional probability bounded away from zero. The next two propositions supply these bounds before we combine them.

Put $n=\lfloor(R-2)/32\rfloor$, $B_k=5kw$, $v_k=1024\mu n(1+\log n)$, and $q=\exp\{-\lambda(56+12\mu)\}$. These quantities are deterministic. The first proposition bounds all routes in the tube through separating rectangles and retains the complete residual line process.

\begin{proposition}\label{prop:upper}
For all sufficiently large admissible $k$, on $E_k$,
\begin{equation}\label{eq:upper-explicit}
\begin{aligned}
 \P(C_k\mid\mathcal G_k)
 &\le \exp\{-qn/\sqrt k\}
 +\exp\Big\{-\frac{n(\log k)^2}{2\lambda B_k}
          +\frac{v_k}{B_k^2}(k-1-\log k)\Big\}\\
 &\le2\exp\Big\{-\frac{k(\log k)^2}{1280\lambda}\Big\}.
\end{aligned}
\end{equation}
\end{proposition}

The second proposition constructs a passage using additional lines in the residual process while preserving sufficiently many distinct full runs. Its uniformity over the initially selected lines is essential for the joint probability. Let $C_0$ denote a finite positive constant depending on the fixed parameters.

\begin{proposition}\label{prop:lower}
The constant $C_0$ can be chosen so that, for all sufficiently large admissible $k$,
\begin{equation}\label{eq:lower}
 \P(F_k\cap C_k\mid\mathcal G_k)\ge\exp\{-C_0k(\log k)^2\}
 \quad\text{on }E_k.
\end{equation}
\end{proposition}

\begin{proof}[Proof of Theorem~\ref{thm:passage}]
Apply Propositions~\ref{prop:upper} and~\ref{prop:lower} at the fixed positive intensity $\lambda$. Absorb the factor two in~\eqref{eq:upper-explicit} by taking $c=1/(2560\lambda)$ and increasing $k_0$. This proves the conditional bounds. Since $E_k$ is $\mathcal G_k$-measurable, integration gives
\[
 \P(E_k)\exp\{-C_0k(\log k)^2\}
 \le\P(E_k\cap F_k\cap C_k)
 \le\P(E_k)\exp\{-ck(\log k)^2\}.
\]
For the upper bound, use $\P(E_k)\le1$. For the lower bound,~\eqref{eq:selection-cost} gives $\log\P(E_k)=-O(k\log k)$, which can be absorbed into $-Ck(\log k)^2$ by increasing $C$. The same enlarged $C$ applies to the conditional bound, and all assertions hold along the prescribed infinite sequence of $k$.
\end{proof}

\section{Proof of the upper bound}

We first state the geometric input controlling the selected lines and the contribution of any additional line. We then prove the upper bound by exposing the remaining environment and asking whether every separating rectangle receives a Cox point. The Poisson exponential formula controls large total arclength in these rectangles, while conditional independence of point counts controls the complementary event. We prove the geometric input after completing this synthesis.

For $1\le j\le n$, set $u_j=32j$, $z_j=(u_j,g(u_j))$, and
\[
 A_j=[u_j-1,u_j+1]\times[g(u_j)-3,g(u_j)+3].
\]
The rectangles are pairwise disjoint and lie horizontally inside $[0,R]$. Define $f_k(\tau)=\sum_{j=1}^n\mathcal H^1(\tau\cap A_j)$ for a line $\tau$. The next lemma supplies the two deterministic bounds used in the probabilistic proof.

\begin{lemma}\label{lem:geometry}
For all sufficiently large admissible $k$, every selection on $E_k$ has the following properties. The function $g$ is $1$-Lipschitz and satisfies $\sup_{0\le x\le R}|g(x)-F_0(x)|\le1$. Every disk of radius $4$ with center having horizontal coordinate in $[0,R]$ meets at most seven lines of $\mathcal D_k$. Moreover, $f_k(\tau)\le B_k$ for every line $\tau$.
\end{lemma}

\begin{proof}[Proof of Proposition~\ref{prop:upper}]
Condition on $\mathcal G_k$ and a selected configuration in $E_k$. The process $\eta=\cL\setminus\mathscr U_k$ is Poisson with intensity $\nu_\mu$ restricted to the complement of $\mathscr U_k$, by the restriction theorem~\cite[Theorem~5.2]{LP17}. Put $T=\sum_{\tau\in\eta}f_k(\tau)$. Each $A_j$ lies in $B(z_j,4)$, so Lemma~\ref{lem:geometry} bounds the selected arclength in $A_j$ by $7\cdot8=56$. Integration first in line displacement and then in angle gives $\int\mathcal H^1(\tau\cap A_j)\nu_\mu(\dd\tau)=\mu|A_j|=12\mu$. Hence, $\E(T\mid\mathcal G_k)\le12\mu n$.

We also need the second moment of the contribution of one line. A line meeting two radius-$4$ disks whose centers are at distance $d\ge32$ has a normal angle in a set of length at most $32/d$. For each angle the possible displacements have length at most $8$. Since $|z_i-z_j|\ge32|i-j|$, and a line contributes at most $8$ arclength to one rectangle,
\begin{equation}\label{eq:variance}
\begin{aligned}
 \nu_\mu\{\tau:\tau\cap B(z_i,4)\ne\varnothing,                      \tau\cap B(z_j,4)\ne\varnothing\}
 &\le\frac{8\mu}{\pi|i-j|}\quad(i\ne j),\\
 \int f_k(\tau)^2\nu_\mu(\dd\tau)
 &\le512\mu n+\frac{1024\mu}{\pi}
                       \sum_{h=1}^{n-1}\frac{n-h}{h}
 \le v_k.
\end{aligned}
\end{equation}
For the angular assertion, the normal projection of the displacement between the centers must have absolute value at most $8$. Its allowed angular length is $2\arcsin(8/d)\le32/d$. The diagonal term uses the intensity $8\mu$ of lines meeting one radius-$4$ disk. Restricting the process to the complement of $\mathscr U_k$ can only decrease these integrals.

For $s\ge0$ and $0\le f\le B_k$, expansion of the exponential gives $\e^{sf}-1-sf\le(f/B_k)^2(\e^{sB_k}-1-sB_k)$. The Poisson exponential formula~\cite[Exercise~3.4]{LP17}, applied on the finite-intensity set of lines meeting the rectangles, and the choice $s=(\log k)/B_k$ therefore yield
\begin{equation}\label{eq:poisson-tail}
\begin{aligned}
 \log\E\big[\exp\{s(T-\E(T\mid\mathcal G_k))\}\mid\mathcal G_k\big]
 &\le\frac{v_k}{B_k^2}(\e^{sB_k}-1-sB_k),\\
 \P\Big(T>12\mu n+\frac{n\log k}{2\lambda}\,\Big|\,\mathcal G_k\Big)
 &\le\exp\Big\{-\frac{n(\log k)^2}{2\lambda B_k}
               +\frac{v_k}{B_k^2}(k-1-\log k)\Big\}.
\end{aligned}
\end{equation}
The bound $f_k\le B_k$ used here is the final assertion of Lemma~\ref{lem:geometry}.

If $A_j$ contains no Cox point, no disk covers the vertical section of $\mathcal T_k$ at $x=u_j$: every possible center of such a disk belongs to $A_j$. Thus, $C_k$ requires every $A_j$ to contain a point. Conditional on all lines, these counts are independent Poisson variables because the rectangles are disjoint; their boundaries have zero arclength almost surely. Write $m_j$ for the total line arclength in $A_j$. Then $\sum_jm_j\le56n+T$. Concavity of $a\mapsto\log(1-\e^{-\lambda a})$ on $a>0$, with zero values handled by continuity, gives
\begin{equation}\label{eq:point-cost}
\begin{aligned}
 \P(C_k\mid\cL)
 &\le\prod_{j=1}^n(1-\e^{-\lambda m_j})
 \le\exp\big\{-n\exp[-\lambda(56+T/n)]\big\},\\
 \P(C_k\mid\mathcal G_k)
 &\le\exp\{-qn/\sqrt k\}
   +\P\Big(T>12\mu n+\frac{n\log k}{2\lambda}\,\Big|\,\mathcal G_k\Big).
\end{aligned}
\end{equation}
This proves the first inequality in~\eqref{eq:upper-explicit} without imposing any restriction on the route within the tube.

For large $k$, $n\ge k^2w/64$ and $B_k=5kw$. The negative exponent in~\eqref{eq:poisson-tail} has magnitude at least $k(\log k)^2/(640\lambda)$. The ratio of the positive term there to this negative term tends to zero: it is at most $2048\mu\lambda(1+\log n)/(5w(\log k)^2)$. The tail term is therefore at most $\exp\{-k(\log k)^2/(1280\lambda)\}$. Finally, $qn/\sqrt k\ge(qw/64)k^{3/2}$, which exceeds $k(\log k)^2/(1280\lambda)$ for all sufficiently large $k$. These estimates give the second inequality in~\eqref{eq:upper-explicit}, uniformly over the selected configuration.
\end{proof}

\begin{proof}[Proof of Lemma~\ref{lem:geometry}]
All selected slopes have absolute value at most one for large $k$, so $g$ is $1$-Lipschitz. Put $a_0=12\varepsilon t$. On its actual supporting segment, $g_i$ differs from the affine extension of its nominal segment by at most $\varepsilon+\varepsilon(1/2+a_0/L)\le2\varepsilon$. Outside the nominal segment but inside the actual one, the nominal piecewise affine graph and this extension differ by at most $2ba_0=6\varepsilon$. Thus, $|g-F_0|\le8\varepsilon<1$ on $[0,R]$.

Partition the selected lines into the first supporting and base lines, four families consisting of the later supports and bases with positive or negative central slope, and the horizontal generators. Within each of the four later families, different centers are separated by at least $2L$. At any horizontal coordinate in $[0,R]$, the errors from the central affine functions are at most $\varepsilon(1+k)$. Consequently, two lines in the same family have vertical separation at least
\[
 4H-2-2\varepsilon(1+k)\longrightarrow\infty.
\]
The term $2$ allows the change from central slope $b$ to $b-1/L$, or from $-b$ to $-b+1/L$, in the base families. A line of absolute slope at most one meeting a radius-$4$ disk has height within $4\sqrt2$ of the center height at the center's horizontal coordinate. Thus, at most one line from each family can meet the disk. The horizontal generators have vertical separations at least $16-4\varepsilon>8\sqrt2$ throughout the range and also contribute at most one. The two first lines contribute at most two, giving seven in total.

Every point $(x,y)$ of an $A_j$ satisfies $|y-F_0(x)|\le5$, using its vertical distance at most $3$ from $g(u_j)$, $|x-u_j|\le1$, and the preceding bounds on $g$. It remains to estimate a line's arclength in this wider tube. A vertical line contributes at most $10$. For a graph line $h(x)=a+vx$ with $|v|\ge2b$, each nominal piece contributes horizontal length at most $20/|v|$. Multiplication by $\sqrt{1+v^2}$ bounds the total by $k(20+10/b)\le40k/b$.

Suppose now that $|v|<2b$, so the conversion from horizontal length to arclength costs less than two. If both $|v-b|$ and $|v+b|$ are at least $b/2$, each piece after the first contributes horizontal length at most $20/b$, and the first contributes at most $L$. If $|v-b|<b/2$, a positive-slope piece with center $c_i$ can be met only when $|h(c_i)|\le H/2+5$. Indeed, at a point of intersection with the wider tube, $|h(c_i)|\le |v-b||x-c_i|+5$, with $|x-c_i|\le L/2$ and $|v-b|<b/2$. At successive such centers, the values of $h$ differ by at least $2vL>bL=2H$. For $H>10$, at most one positive-slope piece can therefore be met. The negative-slope pieces each contribute at most $20/b$, and the first piece still contributes at most $L$. The case $|v+b|<b/2$ is identical with the signs reversed. In every case,
\[
 \mathcal H^1\{(x,y)\in\tau:0\le x\le R, |y-F_0(x)|\le5\}
 \le4L+\frac{40k}{b}=k(4w+160t)\le5kw=B_k.
\]
The last inequality follows from the assumed lower bound on $w/t$. The rectangles are disjoint and lie in this wider tube, so their total contribution $f_k(\tau)$ satisfies the same bound.
\end{proof}

\section{Proof of the lower bound}

We construct a continuous occupied path by adding a group of nearly parallel lines near each supporting segment. These new lines lie in the same angular family as the active support on its own base segment, so they leave the separating points used in the environmental construction available. The lemma below isolates that preservation of full runs. We then combine it with the probability of selecting the new groups and placing enough Cox points along a continuous graph.

Condition on $\mathcal G_k$ and a configuration in $E_k$. Write $v_i$ for the slope of $g_i$, put $\delta=1/100$ and $p=\e^{-4\mu}$, and define
\begin{equation}\label{eq:additional}
 \mathscr D_i=\{h(x)=a+vx:
       |h(c_i)-g_i(c_i)-1/2|\le\delta, |v-v_i|\le\delta/L\}.
\end{equation}
Let $m=\lceil2(R-3)\rceil$, $J=\lceil16\log(m+2)/\lambda\rceil$, and let $D_k$ be the event that exactly $J$ residual lines belong to each $\mathscr D_i$. Let $\mathcal J_k$ be the information in $\mathcal G_k$ together with these additional lines. Finally, put $a_i=\nu_\mu(\mathscr D_i)$, $a_-=4\mu\delta^2/(3\pi L)$, and $a_+=4\mu\delta^2/(\pi L)$.

The preservation statement concerns the remaining line process and requires no condition on the Cox points.

\begin{lemma}\label{lem:preserve}
For all sufficiently large admissible $k$, the sets $\mathscr D_i$ are pairwise disjoint and lie outside $\mathscr U_k$, and
\[
 \P(F_k\mid\mathcal J_k)\ge p/2
 \quad\text{on }E_k\cap D_k.
\]
\end{lemma}

\begin{proof}[Proof of Proposition~\ref{prop:lower}]
Work conditionally on $\mathcal G_k$ as above. The intensity bounds $a_-\le a_i\le a_+$ follow from $\nu_\mu(\dd a\,\dd v)=(\mu/\pi)(1+v^2)^{-3/2}\dd a\dd v$ and $|v|\le1$. By Lemma~\ref{lem:preserve}, the sets are disjoint and outside $\mathscr U_k$. Consequently, conditional on $\mathcal G_k$, their line counts are independent Poisson variables of means $a_i$.

Every chosen line in group $i$ stays within $2\delta$ of $g_i+1/2$ on the corresponding actual segment, since $|x-c_i|\le L/2+a_0$. No condition is imposed on the residual lines outside these groups.

Take horizontal coordinates $s_j=3/2+j(R-3)/m$, $0\le j\le m$, and squares of side $1/4$ centred at $(s_j,g(s_j)+1/2)$. Choose the supporting index active at $s_j$, using either index at a junction. Each of its $J$ additional lines contributes at least $1/8$ arclength to the square: at its center horizontal coordinate the height error is at most $2\delta$, and over horizontal displacement at most $1/16$ the further error is at most $1/16$. Their sum is less than the square's half-height $1/8$. Thus, conditional on the selected additional lines,
\begin{equation}\label{eq:boxes}
 \P(\text{some square has no point from its chosen group})
 \le(m+1)\exp\{-\lambda J/8\}
 \le\frac{m+1}{(m+2)^2}<\frac12.
\end{equation}
The estimate is a union bound and does not require independence between squares or groups used at different squares.

When all squares receive points, choose one point from each. Consecutive chosen points have distance at most $3\sqrt2/4<2$, because $g$ is $1$-Lipschitz and the horizontal spacing is at most $1/2$. The first and last points lie within distance $5\sqrt2/8<1$ of $(1,g(1)+1/2)$ and $(R-1,g(R-1)+1/2)$, respectively. Joining the chosen points and these endpoints by line segments gives a path in the union of their unit disks.

This path also stays in $\mathcal T_k$. Each chosen point differs vertically from $g$ at its own horizontal coordinate by at most $3/4$. Along a segment between two consecutive points, linear interpolation of the two graph values differs from $g$ by at most half their horizontal separation, which is at most $3/8$. Hence, the vertical difference stays below $9/8<2$. The end segments satisfy the same bound, and every point has horizontal coordinate in $[1,R-1]$. This proves that the point event in~\eqref{eq:boxes} implies $C_k$, irrespective of all other lines and points. Cox points on the original selected lines are unnecessary for this lower bound; restoring them and all remaining points preserves the constructed path.

Conditional on $\mathcal J_k$, the favourable point event uses only Cox points on the additional lines, whereas $F_k$ depends only on the remaining line process. These are independent. Lemma~\ref{lem:preserve} and~\eqref{eq:boxes} therefore give conditional probability at least $p/4$ for their intersection. We may now include the full probability cost of the additional groups. Using $J!\le J^J$,
\begin{equation}\label{eq:lower-cost}
\begin{aligned}
 \P(F_k\cap C_k\mid\mathcal G_k)
 &\ge\frac p4\prod_{i=1}^k\frac{\e^{-a_i}a_i^J}{J!}
 \ge\frac p4\exp\{-ka_+-kJ\log(J/a_-)\},\\
 \log\P(F_k\cap C_k\mid\mathcal G_k)
 &\ge-\Big(\frac{32}{\lambda}+o(1)\Big)k(\log k)^2.
\end{aligned}
\end{equation}
Indeed, $R=k^2w$, $J/\log k\to32/\lambda$, $\log(J/a_-)/\log k\to1$, and $ka_+=O_{\mu,t}(1)$. Every estimate is uniform over the initial selected configuration. Increasing a fixed constant $C_0$ proves~\eqref{eq:lower}.
\end{proof}

\begin{proof}[Proof of Lemma~\ref{lem:preserve}]
Each $\mathscr D_i$ is disjoint from its original supporting set because of the height shift $1/2$. Its slope range is disjoint from the original base ranges near the same nominal slope, since their central shift is $1/L$ and $2\varepsilon+\delta<1$. The first, positive, negative, and horizontal slope ranges remain separated for large $k$. Among sets with the same nominal slope and different indices, the centers are at least $2L$ apart, the absolute slope is at least $b/2$, and the allowed heights differ by at most $1/2+4\varepsilon$. Since $bL=2H\to\infty$, these sets cannot overlap. This proves the asserted disjointness, both among the new sets and from $\mathscr U_k$.

We check the full-run count after conditioning on the additional lines. On $I_i$, the vertical distance between the base and its support is at most $1/2+3\varepsilon+(1+2\varepsilon)a_0/L<1$, where $a_0=12\varepsilon t$. For $i\ge2$, the base slope has absolute value between $b/2$ and $b$ and lies strictly between the slopes of the horizontal generators and the active support. Every additional line from $\mathscr D_i$ lies on the support side of this angle ordering, since its slope differs from the support slope by at most $\delta/L$, while the base slope differs by at least $(1-2\varepsilon)/L$.

For $i\ge2$, the base endpoint heights differ from the nominal values $-H,H$ by at most $1/2+3\varepsilon/2+ba_0<1$. The original generators $\gamma_j$ stay within $2\varepsilon$ of height $16j$, and $16N\le H/4$. At the endpoints, the base-generator height difference has opposite signs and absolute value at least $3H/4-1-2\varepsilon$, which eventually exceeds the strip threshold. Thus, each interval on $\ell_i$ at distance at most $2$ from $\gamma_j$ lies inside $I_i$. Its arclength is $4$ divided by the relative sine and is at least $2/b=8t$, because that sine is at most $2b$. The active support fills the other angular family throughout $I_i$. All these intervals therefore remain in the two-family intersection, regardless of additional lines.

For $4\le i\le k$ and $1\le j<N$, take the point $z_{ij}$ on the base at height $16j+8$. Its horizontal coordinate differs from $c_i$ by at most $L/6$ for large $k$, as follows from $|16j+8|\le H/4$, the central height error at most $\varepsilon$, and base slope of absolute value at least $b-(1+\varepsilon)/L$. The point lies inside $I_i$ by this bound and $|y_i-x_i|\le a_0$. The original horizontal generators are at perpendicular distance at least $(8-2\varepsilon)/\sqrt2>2$. A line from another supporting or base group with index $h\ge2$ has center at horizontal distance at least $5L/6$, center height of absolute value at most $1/2+2\varepsilon$, and absolute slope at least $b/2$. Its height differs from that of $z_{ij}$ by at least $7H/12-1/2-2\varepsilon\ge H/2$. This applies to both original and additional lines. For $h=1$, the center height is close to $H/2$, the slope is at least $b/4$, and $z_{ij}$ is at least $17L/6$ to its right, giving the same bound. All these lines are therefore at perpendicular distance greater than $2$ for large $k$.

The own support and additional lines lie in the other angular family, and the base itself is excluded. Hence, none of the fixed lines in the needed angular family covers $z_{ij}$. Conditional on $\mathcal J_k$, the process on the complement of $\mathscr U_k\cup\bigcup_i\mathscr D_i$ remains Poisson with restricted intensity. Since all lines within distance $2$ of a point have intensity $4\mu$, the conditional probability that the needed family leaves each $z_{ij}$ vacant is at least $p$.

Let $Y$ count these vacant points and put $Q=(k-3)(N-1)$. Each generated interval stays within vertical distance $3$ of its level $16j$, so the vacant points lie between consecutive intervals. Distinct vacant points separate consecutive groups of these intervals on their base line. These groups belong to distinct full components, so $S_k\ge Y$, even when other gaps are filled or components extend beyond the segment. The expectation bound and $0\le Y\le Q$ give
\[
 pQ\le\E(Y\mid\mathcal J_k)
 \le pQ/2+Q\P(Y\ge pQ/2\mid\mathcal J_k),\qquad
 \P(S_k\ge pQ/2\mid\mathcal J_k)\ge p/2.
\]
For large $k$, $k-3\ge k/2$ and $N-1\ge H/128$. Therefore, $pQ/2\ge pwk^2/(4096t)\ge(K+4)k^2$ by the assumed lower bound on $w/t$. This is exactly the event $F_k$, and it proves the lemma without assuming independence between the vacant points.
\end{proof}

\section{Passage using only the selected lines}

The upper bound above retains the full random environment. A simpler estimate shows that the originally selected lines themselves admit exponential suppression of passage at small point intensity. This comparison identifies the extra line groups in the lower bound as a substantive contribution from the restored environment.

Let $C_k^{\mathrm{sel}}$ be the same tube-passage event using only Cox points on $\mathcal D_k$. Put $d_k=\lceil(R-4)/2\rceil$. The following estimate permits all changes between selected lines along the passage.

\begin{proposition}\label{prop:selected}
For $0<\lambda<1/28$ and all sufficiently large admissible $k$,
\[
 \P(C_k^{\mathrm{sel}}\mid\mathcal G_k)
 \le2\sqrt2\lambda(2k+N)(28\lambda)^{d_k}
 \quad\text{on }E_k.
\]
\end{proposition}

\begin{proof}
Conditional on $\mathcal G_k$, the selected centers form a Poisson process with intensity measure $\lambda\sum_{\tau\in\mathcal D_k}\mathcal H^1|_\tau$. By Lemma~\ref{lem:geometry}, every disk of radius $2$ centred in the slab $0\le x\le R$ has intensity at most $28\lambda$. The expected number of centers in $[0,2]\times\R$ is at most $2\sqrt2\lambda(2k+N)$, since all selected slopes have absolute value at most one.

An occupied crossing in $\mathcal T_k$ gives a finite chain of distinct centers with successive distances less than $2$, all having horizontal coordinates in $[0,R]$, starting in $[0,2]\times\R$ and ending in $[R-2,R]\times\R$. To see this, cover its compact continuous path by finitely many of the disks meeting the path and take a simple chain in their intersection graph. Every center of a disk meeting the path belongs to the stated slab. Such a chain has at least $d_k$ edges, so its first $d_k$ edges suffice for an upper bound.

The multivariate Mecke formula~\cite[Theorem~4.4]{LP17}, followed by successive integration of these edges, bounds the expected number of such ordered chains by $2\sqrt2\lambda(2k+N)(28\lambda)^{d_k}$. The probability of at least one chain is at most its expectation. This proves the assertion and gives exponential decay in $R=k^2w$ at every fixed $\lambda<1/28$.
\end{proof}

\section{Consequences and remaining scope}

The Cox-point calculation supplies a complete passage estimate for this prescribed oscillating tube. The original selected lines have uniformly bounded local multiplicity, and at small point intensity they alone give an exponential upper bound in the total length. Restoring the Poisson environment changes the probability scale: additional groups of nearly parallel lines along the supporting segments give the lower bound in Theorem~\ref{thm:passage} while preserving the large full-run count, and the separating-rectangle argument gives the corresponding upper bound. All possible occupied routes inside the tube are included in that argument.

The conditional passage cost has negative logarithm of order $k(\log k)^2$, compared with order $k\log k$ for the original line selection. Thus, sampling points does add a larger cost for this geometry. After both costs are included, the negative logarithm is still sublinear in $M=k^2$, so this construction cannot support an exponential upper bound in $M$ even at arbitrarily small fixed $\lambda$.

The upper bound is summable over $k$ for this prescribed construction. It is not an upper bound for the existence of a crossing in an arbitrary tube selected from the full environment, and the path cannot leave the stated tube. Since the theorem holds for every fixed positive point intensity, it does not by itself distinguish the subcritical and supercritical regimes. Extending the control to the relevant family of possible routes, with their selection cost and multiplicity, remains necessary for a subcritical theorem. The result settles the proposed occupied-passage test and supplies no proof that all other approaches fail.

\bibliographystyle{amsalpha}
\bibliography{references}
\end{document}
